%% Interstitial: Paper -1 -> Paper 0
%% Narrative: the structure forced by incompleteness is the sigma-algebra;
%% now we ask when it supports a probability.

Chapter~\ref{chap:paper-minus-one} established that a coherent observer
cannot avoid the $\sigma$-algebra.  Beginning only with a Boolean algebra
$\mathcal{E}_Q$ of distinguishable events and a finitely-additive
valuation $\ell_Q$, the commitment to acting as if a coherent world
exists forces continuity at $\emptyset$ --- and continuity at
$\emptyset$ is precisely the condition under which $\ell_Q$ extends
uniquely to a $\sigma$-additive measure on $\sigma(\mathcal{E}_Q)$.

The $\sigma$-algebra, then, is not an assumption.  It is the form
that the observer's coherent commitments must take.

\medskip

But having a $\sigma$-algebra on each outcome space $O_Q$ is not yet
having a probability on the system as a whole.  The observer has a
family of local valuations $\{\ell_Q\}$ --- one for each query --- and
these valuations are compatible across refinements.  The next question
is: does this compatible family determine a single global probability
measure on the realization space $\Omega$?

This is Stopping Point~2, and it is a strictly harder question than
SP1.  SP1 concerned a single observer at a single level: can they
extend their finitely-additive valuation to their own $\sigma$-algebra?
SP2 concerns the whole system: can the compatible family be assembled
into a $\sigma$-additive measure on $\Omega = \varprojlim O_Q$?

\medskip

Chapter~\ref{chap:paper-zero} addresses SP2.  It does so by identifying
a condition --- surjective projective uniform tightness (SPUT) --- under
which the compatible family concentrates on $\Omega$ rather than
dissipating into the product $\prod_Q O_Q$.  The condition is
topological: it asks that for every $\varepsilon > 0$, the family
$\{\ell_Q\}$ is supported, up to $\varepsilon$, on a compact
surjectively-compatible subfamily $\{K_Q\}$.  Compactness of the $K_Q$
is the external witness that SP1 showed was necessary: it is the
topological expression of the observer's faith in a world that closes
the gaps.

Two architectural cases are established.  When outcome spaces are
compact Hausdorff, SPUT holds automatically and $\sigma$-additivity
is a structural theorem.  When outcome spaces are standard Borel,
Musia\l's theorem provides a topology-free route via perfectness.
Both cases close SP2 for the architectures that arise in the
subsequent papers of this thesis.

\medskip

The passage from Chapter~\ref{chap:paper-minus-one} to
Chapter~\ref{chap:paper-zero} is therefore the passage from
\emph{local} to \emph{global}: from a single observer's forced
commitments to the assembly of all observers' commitments into a
single coherent probability on the world they share.
